\documentclass[11pt,a4paper]{article}
\usepackage[margin=22mm]{geometry}
\usepackage[T1]{fontenc}
\usepackage{lmodern}
\usepackage{amsmath,booktabs,tabularx,array}
\usepackage[hidelinks]{hyperref}
\usepackage{enumitem}
\usepackage{fancyhdr}
\setlength{\parindent}{0pt}
\setlength{\parskip}{6pt}
\setlist[itemize]{nosep,leftmargin=*}
\renewcommand{\arraystretch}{1.15}
\pagestyle{fancy}
\fancyhf{}
\fancyhead[L]{Hong Kong IPO research}
\fancyhead[R]{Progress report}
\fancyfoot[C]{\thepage}
\setlength{\headheight}{14pt}
\newcommand{\src}[1]{\par{\footnotesize\textit{Source:} \nolinkurl{#1}}}
\newcommand{\pct}{\%}
\title{Hong Kong IPO Research\\Research, Data, and Results}
\author{Research progress report}
\date{Report date: 2 October 2026\\Data cutoff: 30 September 2026}
\begin{document}
\maketitle

\section{Research purpose}
This research examines initial public offerings (IPOs) on the Hong Kong Main Board.
The sample contains ordinary IPOs listed in 2026.
The main question is how offer prices, investor demand, and share allocation relate to investor returns.

The research examines five related questions.
\begin{itemize}
\item How large are first-day returns, and how do they change across listing periods?
\item Do high first-day returns give retail applicants high returns on application funds?
\item How do cornerstone allocations and retail demand relate to first-day returns?
\item Do existing A-share prices help explain H-share offer prices and later prices?
\item What happens to returns after listing, lockup expiry, and the end of stabilization?
\end{itemize}

An IPO can have a high first-day return and a low return on application funds.
An applicant receives only part of the shares requested.
A high return on allocated shares does not apply to all application funds.
This difference is central to the retail analysis.

The results also test possible information and supply channels.
Cornerstone investors can give information about an issuer.
Their locked shares can also reduce the shares available for trade.
An existing A-share price gives a public price reference.
These channels can operate at the same time.

\section{Current research state}
The project has a data pipeline and reproducible analysis scripts.
The current sample contains 113 issuers.
The analysis covers listings from 2 January to 30 September 2026.
The data cutoff is 30 September 2026.
The report date does not extend the data cutoff.

The project has completed price corrections, selected field repairs, and contract-based lockup records.
It has also completed descriptive tables, regressions, retail-return calculations, and event studies.
The results are exploratory.
The project has already examined the data and tested several models.
These tests are not preregistered tests.

The results describe associations.
They do not establish that investor demand, cornerstone participation, or a listing route causes a return change.

{\small\textit{Language note.} This report uses short sentences and defined technical names, with ASD-STE100 as its writing guide.
It is not a certified dictionary-compliance review.
The official guide is available at \url{https://www.asd-ste100.org/}.}

\newpage
\section{Data and completed repairs}
\subsection{Sample and sources}
\begin{table}[h]
\centering
\caption{Current sample and selected data coverage}
\begin{tabular}{lr}
\toprule
Data group & Issuers or records\\
\midrule
First quarter listings & 38\\
Second quarter listings & 45\\
Third quarter listings & 30\\
All 2026 listings in the study & 113\\
A+H issuers, independent flag & 38\\
Main regression, complete cases & 105\\
Issuers with margin-financing reports & 25\\
Source-reported margin records & 93\\
\bottomrule
\end{tabular}
\end{table}

Official listing records define issuer membership.
Prospectuses supply issuer, offer, financial, investor, and contract information.
Allotment announcements supply final allocation and subscription information.
Market data supply stock prices, trading volumes, and index prices.
Dated media reports supply margin-financing observations.
These media observations do not replace official subscription outcomes.

The pipeline stores source references, extracted values, validation records, and review records.
It exports cohort data and builds the master CSV.
The analysis selects 2026 issuers by their actual listing dates.
Historical issuers can remain in storage without entering this study.
Unknown values remain missing.
Each model reports its own usable sample.

\subsection{Repairs already completed}
\begin{itemize}
\item PR \#20 corrected raw first-day price inputs for 26 issuers in the original 106-issuer sample.
\item PR \#23 applied independent field reviews to 14 second-quarter board-seat records.
\item These records now contain eight yes values, one no value, and five unknown values.
\item PR \#23 defined the units for 6228.HK: one Hong Kong depositary receipt (HDR) represents ten underlying shares.
\item The retail calculation now includes 6228.HK with the corrected unit treatment.
\item PR \#23 and PR \#25 added contract-based lockup evidence for the 113-issuer sample.
\item PR \#25 added seven third-quarter issuers and changed the total sample from 106 to 113.
\item PR \#27 completed the stored table directories for all 113 issuers.
\end{itemize}

A review of one field does not verify all fields for that issuer.
A stored source quote does not, by itself, establish correct interpretation.
A complete price window does not establish an external cause for a price change.

\src{pipeline/exports/HKIPO-MB-MASTER_clean.csv}
\src{pipeline/reports/research_readiness/README.md}

\newpage
\section{Measures and analysis methods}
\subsection{Return measures}
Let $P_{0i}$ be issuer $i$'s offer price.
Let $P_{1i}$ be its raw first-day closing price.
The first-day return is
\begin{equation}
IR_i=\frac{P_{1i}}{P_{0i}}-1.
\end{equation}
The main regression uses $y_i=\ln(1+IR_i)$.
This transformation reduces the influence of very large positive returns.
A coefficient on this outcome is not a percentage-point change in $IR$.

Let $a_i$ be final public-offer units divided by valid applied units.
It is an issuer-level average allocation rate.
The retail analysis reports three different measures:
\begin{align}
\overline{IR}&=\frac{1}{N}\sum_i IR_i,\\
IR_{\mathrm{allocation}}&=\frac{\sum_i a_iIR_i}{\sum_i a_i},\\
R_{\mathrm{application}}&=\frac{1}{N}\sum_i a_iIR_i.
\end{align}
The first measure gives equal weight to each IPO.
The second measure weights first-day returns by allocation rates.
The third measure assumes equal application funds for each IPO.
It measures gross return on those application funds.
It does not give an individual applicant's allocation probability.

An aftermarket buy-and-hold return starts at the first-day close.
It excludes the first-day return.
The A+H price gap compares the H price with the same issuer's A price in Hong Kong dollars.
The exchange-rate measure is Hong Kong dollars per Chinese yuan.
An event CAR adds daily stock returns less benchmark returns over an event window.

\subsection{Regression and inference}
The main baseline model has the following form:
\begin{equation}
\ln(1+IR_i)=\alpha+X_i'\beta+\varepsilon_i.
\end{equation}
The vector $X_i$ includes age, offer size, A+H status, investor backing, sponsor quality, and cornerstone allocation.
It also includes prior market returns, prior IPO count, and an April--June indicator.
A fourth model adds the final retail subscription ratio.
All four nested models use the same 105 complete cases.

The analysis reports HC3 standard errors.
HC3 permits unequal error variance across issuers.
The analysis also groups errors by listing month.
There are only nine listing months.
The restricted wild cluster bootstrap checks inference with these few groups.
It evaluates all 512 sign combinations for the baseline model.

The April--June indicator came from observed data patterns.
It is an exploratory time control.
Final subscription demand can respond to expected returns.
The model with final demand does not identify a pricing cause.
\src{analysis/out/module_b/sample_selection.md}
\src{analysis/out/module_b/table6_inference.md}

\newpage
\section{First-day returns and regression results}
\subsection{Descriptive results}
\begin{table}[h]
\centering
\caption{First-day returns by listing quarter}
\begin{tabular}{lrrrr}
\toprule
Measure & Q1 & Q2 & Q3 & All\\
\midrule
Issuers & 38 & 45 & 30 & 113\\
Mean return (\%) & 33.9 & 92.7 & 18.3 & 53.2\\
Median return (\%) & 13.4 & 80.0 & 0.0 & 14.8\\
Negative returns (\% of issuers) & 10.5 & 17.8 & 46.7 & 23.0\\
\bottomrule
\end{tabular}
\end{table}

The full-sample mean is much larger than the median.
Large positive returns raise the mean.
The second quarter has the highest mean and median.
Almost half of third-quarter issuers have negative first-day returns.
These differences show substantial variation within 2026.
They do not identify the cause of that variation.
\src{analysis/out/module_a/table1_stylized_facts.md}

\subsection{Main baseline model}
\begin{table}[h]
\centering
\caption{Selected baseline coefficients: 105 issuers, nine listing months}
\begin{tabular}{lrrr}
\toprule
Variable & Coefficient & HC3 $p$ & Wild cluster $p$\\
\midrule
A+H issuer & $-0.313$ & 0.018 & 0.004\\
Cornerstone allocation & $-0.353$ & 0.368 & 0.270\\
VC/PE backing & 0.088 & 0.577 & 0.770\\
Top-tier sponsor & 0.051 & 0.612 & 0.555\\
Prior 20-day HSI return & 2.055 & 0.033 & 0.020\\
April--June indicator & 0.394 & 0.039 & 0.039\\
\bottomrule
\end{tabular}
\end{table}

A+H issuers have lower first-day price ratios after the included controls.
The coefficient implies approximately $\exp(-0.313)-1=-26.9\%$ for the price ratio $1+IR$.
It does not imply a 26.9 percentage-point decrease in first-day return.
A+H status can also reflect issuer selection and firm differences.

The baseline does not give precise evidence for cornerstone allocation, investor backing, or sponsor quality.
These results do not prove that their effects are zero.
The model has a small sample and several explanatory variables.

The fourth model's subscription coefficient is 0.095, with an HC3 standard error of 0.032.
Its adjusted $R^2$ is 0.264, compared with 0.194 for the baseline.
Final demand improves the fit within the sample.
This improvement does not resolve reverse causality.
\src{analysis/out/module_b/table4_regressions.md}
\src{analysis/out/module_b/table6_inference.md}

\newpage
\section{Retail application returns and cornerstone results}
\subsection{Returns have different denominators}
\begin{table}[h]
\centering
\caption{Retail measures: 113 issuers}
\begin{tabular}{lrrr}
\toprule
Measure (\%) & Estimate & Lower bound & Upper bound\\
\midrule
Equal-weight first-day return & 53.153 & 28.904 & 77.053\\
Allocation-weighted return & 1.165 & $-2.239$ & 8.574\\
Return on application funds & 0.020 & $-0.051$ & 0.098\\
\bottomrule
\end{tabular}
\end{table}

The bounds come from 2,000 bootstrap samples drawn by listing month.
The procedure uses nine months and the fixed seed 20260930.
These bounds are descriptive because there are few months.

A HK\$10,000 application to each IPO gives an average gross profit of approximately HK\$1.97 under the assumed allocation rates.
High first-day returns therefore do not imply high returns on application funds.
The allocation rate changes the amount exposed to the first-day price increase.

The calculation uses an average issuer allocation rate.
Actual allotment can differ by application size.
The calculation does not model a specific application tier or an individual account.
The financing calculations use assumed fees, borrowing shares, rates, and a two-day holding period.
They are cost scenarios, not observed borrowing outcomes.
The calculation excludes some transaction costs and opportunity costs.
It is not a complete net-return measure.
\src{analysis/out/research_frontier/retail_summary.csv}
\src{analysis/out/research_frontier/research.md}

\subsection{Cornerstone sensitivity analysis}
\begin{table}[h]
\centering
\caption{Cornerstone allocation: common sample of 106 issuers}
\begin{tabular}{lrrr}
\toprule
Model & Coefficient & HC3 $p$ & Wild cluster $p$\\
\midrule
Controls known before final demand & $-0.309$ & 0.451 & 0.191\\
Controls plus final demand & $-0.455$ & 0.213 & 0.082\\
\bottomrule
\end{tabular}
\end{table}

Both coefficients are negative, but the inference depends on the method.
The analysis also measures sensitivity to omitted variables.
This check describes how an unobserved factor can change the point estimate.
It does not supply a causal confidence interval.

Cornerstone shares and demand are issuer and investor choices.
Adding demand can control for part of the channel under study.
It can also introduce selection bias.
The current results do not separate certification from a reduction in tradable supply.
\src{analysis/out/research_frontier/research.md}

\newpage
\section{A+H prices, margin financing, and later returns}
\subsection{A+H price reference}
The independent A+H sample contains 38 issuers.
The mean H offer-price gap is $-38.2\%$ relative to the A-share reference price.
The mean first-day closing gap is $-30.7\%$.
The mean first-day return for these issuers is 9.4\%.
The two price-gap measures use different reference dates.
Their difference also contains A-share and exchange-rate changes.

Some A+H models use 34 valid observations.
At trading day 60, only 25 issuers have matched price pairs.
Their mean gap is $-28.7\%$ at both endpoints.
The paired test has $p=0.996$.
This result does not prove that price convergence is absent.

The market-model mean A-share CAR is $-7.6\%$ around the H listing date, for the $[-5,+5]$ window.
This estimate uses 34 issuers; its placebo $p$-value is 0.004.
Other announcements, industry movements, and listing choices can affect this result.
The current A+H report's rank-correlation entry is unavailable.
This report does not interpret that entry as a numerical result.
\src{analysis/out/ah_anchor/ah_anchor.md}

\subsection{Margin financing}
The dataset has 93 reported observations for 25 of the 113 issuers.
The sample depends on media survey coverage.
It is not a random sample of all IPOs.

The model with the latest demonstrably public pre-deadline observation uses 25 issuers and six listing months.
After time-window and size controls, its coefficient is 0.072.
The HC3 $p$-value is 0.384; the wild cluster $p$-value is 0.531.
The current evidence does not give a precise controlled association.
Publication before subscription close does not establish publication before pricing.
\src{analysis/out/margin/margin.md}

\subsection{Returns after listing and around supply events}
\begin{table}[h]
\centering
\caption{Buy-and-hold price returns from the first-day close}
\begin{tabular}{lrrr}
\toprule
Horizon & Issuers & Mean (\%) & Median (\%)\\
\midrule
Day 20 & 102 & 5.5 & $-1.0$\\
Three months & 82 & 6.6 & $-16.6$\\
Six months & 37 & 29.7 & $-22.3$\\
\bottomrule
\end{tabular}
\end{table}

The six-month sample contains first-quarter listings only.
The positive mean and negative median show a large influence from a few positive returns.
Different horizons contain different issuers.
They cannot represent one balanced return path for all 113 IPOs.

The lockup $[-5,+5]$ window has 29 issuers and a mean HSI-adjusted CAR of $-5.6\%$.
Its placebo $p$-value is 0.100.
The stabilization-end window has 31 issuers and a mean CAR of $-2.2\%$.
Its placebo $p$-value is 0.524.
These findings do not establish a general causal supply effect.
\src{analysis/out/module_a/table2_aftermarket.md}
\src{analysis/out/academic/events.md}

\newpage
\section{Other completed tests and remaining work}
\subsection{Other completed analysis}
The partial-adjustment analysis uses 43 range-priced offers.
It tests whether price revision relates to the first-day return.
The controlled revision coefficient is 1.52; its wild cluster $p$-value is 0.262.
The estimate is too imprecise for a strong conclusion.
Fixed-price offers do not have a defined position within a price range.

The project also tests sponsor differences, underwriting fees, market crowding, and extreme returns.
Sponsor and fee analyses use documented proxies and selected source fields.
Some generated reports retain old explanatory notes.
Those notes require a separate consistency review before publication.

Prediction tests hold out issuers or listing months.
The prediction sample has 90 issuers with the required prior-return input.
The baseline leave-one-month-out $R^2$ is 0.150.
The April--June-only model gives 0.142.
These tests use a time pattern selected after observing the data.
They do not constitute a real-time forecast of a future high-return period.
\src{analysis/out/academic/partial_adjustment.md}
\src{analysis/out/extended/prediction.md}

\subsection{Limits that remain}
\begin{itemize}
\item The study covers nine months of one year.
\item Some fields remain unknown, and some return windows are incomplete.
\item A few large returns strongly affect mean results.
\item Listing choices, investor participation, and demand can share unobserved causes.
\item Mechanism groups have limited comparable observations.
\item All 19 Chapter 18C issuers use Mechanism A in the current support table.
\item Many models and tests have already been examined.
\item Statistical sensitivity checks do not create external treatment variation.
\end{itemize}

\subsection{Next research steps}
\begin{enumerate}[nosep,leftmargin=*]
\item Select one main question and define its economic mechanism.
\item Fix the main outcome, sample rules, and inference method before additional tests.
\item For retail returns, collect official allotment tables by application size.
\item For A+H prices, collect first issuance announcement dates and comparable firms.
\item Add complete event windows as they become available.
\item Freeze a prediction model before testing it on fourth-quarter listings.
\item Review source-to-table links and remaining old report notes before paper submission.
\end{enumerate}

\section{Technical names and result sources}
\begin{tabularx}{\textwidth}{@{}lX@{}}
\toprule
Name & Meaning in this report\\
\midrule
IPO & Initial public offering.\\
Issuer & The company that offers securities.\\
A+H issuer & A company with A shares and listed H shares.\\
Cornerstone investor & An investor with an agreed IPO allocation and contractual restrictions.\\
VC/PE & Venture capital or private equity investor backing.\\
HSI & Hang Seng Index, used as a market benchmark.\\
CAR & Cumulative abnormal return over the specified event window.\\
Lockup & A contractual restriction on disposal of securities.\\
Stabilization & Permitted price-support activity after an offering.\\
\bottomrule
\end{tabularx}

The source notes identify the stored outputs used for this report.
The report explains existing results; it does not rerun the analysis.
The research entry is \nolinkurl{docs/RESEARCH_START.md}.
The current plan is \nolinkurl{docs/RESEARCH_PLAN_2026.md}.
The repair record is \nolinkurl{pipeline/reports/research_readiness/README.md}.
\end{document}
